%%% FIELD proof-hall-coverage
By \(\text{ass:ordinal-support}\), \(L\geq 3\), so the declared support \(X=\{0,\ldots,L-1\}\) is finite and nonempty.  Put \(p(i)=\mu_{1,i}\) and \(q(j)=\mu_{0,j}\).  Since \(\mu_1,\mu_0\in\Delta_L\), \(p\) and \(q\) are probability measures on \(X\).  For the fixed set \(S\subseteq E\), define the outside relation
\[
R_S=\{(i,j)\in X\times X:i-j\notin S\}.
\]
For every \(A\subseteq X\), the neighborhood of \(A\) under \(R_S\) is
\[
N_{R_S}(A)=\{j\in X:\exists i\in A\text{ such that }i-j\notin S\}
=N_{S^{\mathsf c}}(A),
\]
where the last equality is \(\text{def:complement-neighborhood}\).

First, the optimization in \(\text{def:worst-coverage}\) is well posed at boundary marginal laws.  Indeed, the product table \(\gamma^{\otimes}_{ij}=p(i)q(j)\) is nonnegative and satisfies
\[
\sum_{j\in X}\gamma^{\otimes}_{ij}=p(i)\sum_{j\in X}q(j)=p(i),
\qquad
\sum_{i\in X}\gamma^{\otimes}_{ij}=q(j)\sum_{i\in X}p(i)=q(j).
\]
Thus \(\gamma^{\otimes}\in\Gamma(\mu_1,\mu_0)\) by \(\text{def:coupling-polytope}\), even when some coordinates of either marginal are zero.  In particular, the feasible coupling set is nonempty.  Attainment of the stated minimum will follow below from the coupling furnished by \(\text{lem:finite-strassen-deficiency}\), so no positivity or interiority condition is needed.

Define the finite Hall deficiency
\[
\delta_S=\max_{A\subseteq X}\bigl\{p(A)-q(N_{S^{\mathsf c}}(A))\bigr\}.
\]
The maximum is attained because \(2^X\) is finite.  Its term at \(A=\varnothing\) is zero, while every term is at most \(p(A)\leq1\); consequently
\[
0\leq\delta_S\leq1.
\]
For an arbitrary \(\gamma\in\Gamma(\mu_1,\mu_0)\), let
\[
e_S(\gamma)=\sum_{\substack{i,j\in X\\i-j\in S}}\gamma_{ij}.
\]
The sets \(R_S\) and \(\{(i,j):i-j\in S\}\) partition \(X\times X\), and every coupling has total mass one.  Therefore
\[
\gamma(R_S)=1-e_S(\gamma).
\]
Apply the necessity direction of \(\text{lem:finite-strassen-deficiency}\) to the finite sets \(U=V=X\), relation \(R_S\), marginals \(p,q\), and the nonnegative number \(\varepsilon=e_S(\gamma)\).  It yields, for every \(A\subseteq X\),
\[
p(A)\leq q(N_{R_S}(A))+e_S(\gamma)
=q(N_{S^{\mathsf c}}(A))+e_S(\gamma).
\]
Taking the maximum over \(A\) gives
\[
\delta_S\leq e_S(\gamma).
\]
Since this holds for every feasible coupling, \(\delta_S\leq h_\mu(S)\) by \(\text{def:worst-coverage}\).

Conversely, the definition of \(\delta_S\) gives the simultaneous inequalities
\[
p(A)\leq q(N_{R_S}(A))+\delta_S
\qquad\text{for every }A\subseteq X.
\]
Because \(\delta_S\geq0\), the sufficiency direction of \(\text{lem:finite-strassen-deficiency}\), with \(\varepsilon=\delta_S\), supplies a coupling \(\widetilde\gamma\in\Gamma(\mu_1,\mu_0)\) satisfying
\[
\widetilde\gamma(R_S)\geq1-\delta_S.
\]
Taking complementary mass gives
\[
e_S(\widetilde\gamma)=1-\widetilde\gamma(R_S)\leq\delta_S.
\]
Thus \(h_\mu(S)\leq e_S(\widetilde\gamma)\leq\delta_S\).  Combining this with the preceding lower bound proves
\[
h_\mu(S)=\delta_S
=\max_{A\subseteq X}\{\mu_1(A)-\mu_0(N_{S^{\mathsf c}}(A))\},
\]
and also shows directly that \(\widetilde\gamma\) attains the minimum.

It remains to convert the Hall deficiency into the rectangle representation.  To make the boundary case unambiguous, read the rightmost displayed maximum in the statement as the maximum of the finite nonempty set
\[
\mathcal T_S=\{0\}\cup
\{p(A)+q(B)-1:\varnothing\neq A,B\subseteq X,\ A-B\subseteq S\}.
\]
This agrees with the displayed nested notation whenever a nonempty feasible rectangle exists and remains defined when none exists.  For each \(A\subseteq X\), set
\[
B_A=X\setminus N_{S^{\mathsf c}}(A).
\]
If \(i\in A\) and \(j\in B_A\), the definition of \(B_A\) says that no element of \(A\), in particular \(i\), has \(i-j\notin S\).  Hence \(i-j\in S\), and \(\text{def:difference-set}\) gives
\[
A-B_A\subseteq S.
\]
Since \(B_A\) is the complement of \(N_{S^{\mathsf c}}(A)\) in \(X\) and \(q(X)=1\),
\[
p(A)-q(N_{S^{\mathsf c}}(A))
=p(A)+q(B_A)-1.
\]
Choose \(A^\star\subseteq X\) attaining \(\delta_S\).  If \(\delta_S>0\), then \(A^\star\neq\varnothing\), because the deficiency at the empty set is zero.  Also \(B_{A^\star}\neq\varnothing\), because otherwise
\[
p(A^\star)+q(B_{A^\star})-1=p(A^\star)-1\leq0,
\]
contrary to \(\delta_S>0\).  Therefore, when \(\delta_S>0\), the value \(\delta_S\) is one of the rectangle values in \(\mathcal T_S\).  When \(\delta_S=0\), it is the explicit element zero.  In either case,
\[
\delta_S\leq\max\mathcal T_S.
\]

For the reverse inequality, take nonempty \(A,B\subseteq X\) with \(A-B\subseteq S\).  If \(j\in B\), then \(i-j\in S\) for every \(i\in A\), by \(\text{def:difference-set}\).  Thus \(j\notin N_{S^{\mathsf c}}(A)\), whence \(B\subseteq B_A\).  Nonnegativity of the probability measure \(q\) now gives
\[
\begin{aligned}
p(A)+q(B)-1
&\leq p(A)+q(B_A)-1\\
&=p(A)-q(N_{S^{\mathsf c}}(A))\\
&\leq\delta_S.
\end{aligned}
\]
The remaining member zero of \(\mathcal T_S\) is also at most \(\delta_S\).  Hence \(\max\mathcal T_S\leq\delta_S\), and consequently
\[
h_\mu(S)
=\max_{A\subseteq X}\{\mu_1(A)-\mu_0(N_{S^{\mathsf c}}(A))\}
=\max\!\left\{0,
\max_{\substack{\varnothing\neq A,B\subseteq X\\A-B\subseteq S}}
\bigl(\mu_1(A)+\mu_0(B)-1\bigr)\right\}.
\]
The Hall maximum and the maximum over \(\mathcal T_S\) are maxima of finite nonempty sets and therefore are attained.  When their common value is positive, the construction above supplies an attaining nonempty rectangle; when it is zero and no such rectangle exists, the explicit zero supplies the maximizer.  The coupling \(\widetilde\gamma\) constructed above attains the coupling minimum.  This proves every assertion without imposing positivity of any marginal cell.
